\chapter{Results}

ODSA produces four results that advance the characterisation of the Antarctic bed. First,
neither $\beta$ nor relief separates any pair of the published landscape classes~\cite{Ockenden_2026} on co-located windows, and no element of the vector separates more than half of the class pairs, because the distributions of each element overlap between classes. Landscape classification from radar profiles therefore requires several elements measured together. Second, a radar profile constrains the landscape class without elements measured together. Second, a radar profile constrains the landscape class without identifying it uniquely: of the 11 archetypes in the catalogue, a median of two remain admissible per window over 651 windows. Third, the bed texture elements largely decorrelate within one 50~km window length along a profile, while bed elevation and surface speed remain correlated over hundreds of kilometres, and the bed class of windows on neighbouring flight lines 1 to 25~km apart agrees only slightly more often than expected by chance. Bed texture therefore changes between  neighbouring flight lines and cannot be recovered by interpolating between them. Fourth, in every region the detection limit for a difference in $\beta$ between flow-parallel and flow-perpendicular profiles is two to twenty times larger than the difference measured on deglaciated beds with mapped glacial fabric. Reducing this limit requires one to two orders of magnitude more coverage from flight lines that cross the same ground at several angles.


% \begin{figure}
%     \includegraphics[scale=0.37]{figures/catalogue_examples.png}
%     \caption{} 
%     \label{fig:catalogue_examples}
% \end{figure}


